% =====================================================================
% CHAPTER 3 — LITERATURE REVIEW (draft opening for review)
% Intended placement: NEW Section 3.1, slotted BEFORE the current
% conceptual-framework sections (which become 3.2 reactive/proactive,
% 3.3 stakeholder pressure, 3.4 the bridge, 3.5 the playing field).
% Built only from sources already in references.tex.
% %CL = notes for you; VERIFY each citation's fit before locking (per our source rule).
% Missing from references.tex and needed: Mitchell, Agle & Wood (1997) — cited in 3.3/3.5.
% =====================================================================

\subsection{3.1 Literature review}\label{literature-review}

This section situates the thesis within five strands of existing work: the
growth and fragmentation of corporate sustainability reporting; the particular
position of SMEs within it; the reactive--proactive posture and maturity
literature; the theoretical lenses used to read multi-stakeholder demand; and the
governance of seafood sustainability specifically. It closes by identifying the
gap the thesis addresses.

\textbf{The growth and fragmentation of sustainability reporting.} Corporate
sustainability reporting has moved over two decades from a voluntary,
largely narrative practice to an increasingly regulated and, in the European
Union, mandatory field (Hummel \& Jobst, 2024; Dinh, Husmann \& Melloni, 2023).
Its centre of gravity has shifted toward standardised, decision-useful metrics
and the double-materiality principle (Nielsen, 2023; Nahrstedt, 2026). Yet the
same period produced a proliferation of overlapping standards, frameworks and
ratings that observers have repeatedly criticised as fragmented and
insufficiently comparable, prompting sustained calls for harmonisation (Afolabi,
Ram \& Rimmel, 2022). A growing body of work now examines the European
Sustainability Reporting Standards (ESRS) directly --- their information
architecture (Fragidis \& Papafloratos, 2025), their operationalisation into
assessment frameworks (Ahmad et al., 2025), and their implications for
manufacturing firms (Bataleblu et al., 2024) --- but this literature is
overwhelmingly oriented toward the large undertakings the standards were written
for.

\textbf{SMEs and the proportionality problem.} A distinct strand documents the
gap between this expanding reporting apparatus and the capacity of small and
medium-sized enterprises. SMEs operate with limited financial and managerial
resources and little dedicated reporting expertise, which constrains their
environmental and sustainability engagement (Battisti \& Perry, 2011). The
proportionality response --- reduced-scope voluntary standards such as the VSME
--- has begun to attract study, including of its applicability to agri-food SMEs
specifically (Tugliani, Guareschi \& Arfini, n.d.). What this strand establishes
is that lighter standards reduce the \emph{volume} of what an SME must report, but
do not resolve the prior question of \emph{where}, within a firm's own mix of
demands, a proportionate effort should be focused.

\textbf{Reactive and proactive postures, and reporting maturity.} A long line of
strategy research distinguishes reactive from proactive environmental and
sustainability postures and links the proactive posture to the development of
valuable organisational capabilities and to long-run performance (Sharma \&
Vredenburg, 1998; Torugsa, O'Donohue \& Hecker, 2013; Chan, Lai \& Kim, 2022;
Fazli et al., 2023). A parallel maturity literature characterises the stages
through which firms progress and the difficulty of moving between them
(Baumgartner \& Ebner, 2010; Cöster et al., 2020; Lozano, 2013). Together these
provide the posture-and-maturity vocabulary the thesis adopts, but they are
largely pitched at the level of the individual firm's strategy and say little
about how the \emph{external structure of demand} shapes whether a small firm can
make the reactive-to-proactive transition at all.

\textbf{Theoretical lenses on multi-stakeholder demand.} Three established lenses
help read that external structure. Stakeholder-salience theory (Mitchell, Agle \&
Wood, 1997) explains how firms triage competing claims by power, legitimacy and
urgency. Institutional theory (DiMaggio \& Powell, 1983) explains why disclosure
conventions converge across firms through coercive, mimetic and normative
isomorphism even where formally voluntary. Legitimacy theory (Suchman, 1995)
distinguishes pragmatic, moral and cognitive legitimacy, a gradient that maps
onto the progression from legally enforced to normatively expected obligation.
The thesis draws these together in its conceptual framework (Sections~3.2--3.4).

\textbf{The governance of seafood sustainability.} Finally, a sector-specific
literature characterises seafood sustainability governance as a
crowded, multi-actor field in which regulators, retailers, certifiers, NGOs and
investors exert parallel and overlapping demands (Bush \& Oosterveer, 2019;
Barclay \& Miller, 2018), coordinated through certification schemes, improvement
projects and voluntary initiatives. This work describes the governance landscape richly, but treats it largely from the
standpoint of the movement and its institutions rather than from the standpoint
of a single resource-limited firm trying to answer all of them at once.

\textbf{The gap.} Across these strands, the reporting literature maps frameworks,
the strategy literature studies posture, the theoretical literature explains
multi-stakeholder pressure, and the seafood literature describes sector
governance --- but they remain largely separate. No existing work systematically
maps cross-stakeholder disclosure \emph{demand} at the level of a single sector
and firm, in a way that tells a resource-limited SME which disclosure topics, if
built first, would satisfy the largest share of the demands it actually faces.
Supplying that mapping --- and the diagnostic that follows from it --- is the
contribution this thesis makes. The remainder of the chapter builds the
conceptual framework through which it is done.
